% =========================================================
% ppt/src/minggu-15.tex
% Slide pertemuan 15
% Bab XVI — Debugging Level Instruksi
%
% Sub-CPMK 4.1--4.3 | CPMK04 | CPL03
% =========================================================

\documentclass[aspectratio=169,11pt]{beamer}
% =========================================================
% ppt/src/preamble-ppt.tex
% Tema, warna, paket, dan makro bersama seluruh deck slide.
%
% Dipakai oleh setiap berkas minggu-NN.tex dengan pola:
%   \documentclass[aspectratio=169,11pt]{beamer}
%   % =========================================================
% ppt/src/preamble-ppt.tex
% Tema, warna, paket, dan makro bersama seluruh deck slide.
%
% Dipakai oleh setiap berkas minggu-NN.tex dengan pola:
%   \documentclass[aspectratio=169,11pt]{beamer}
%   % =========================================================
% ppt/src/preamble-ppt.tex
% Tema, warna, paket, dan makro bersama seluruh deck slide.
%
% Dipakai oleh setiap berkas minggu-NN.tex dengan pola:
%   \documentclass[aspectratio=169,11pt]{beamer}
%   \input{preamble-ppt}
%
% Palet warna disalin dari buku_ajar/config/styles.tex agar
% identitas visual slide sama dengan buku ajar.
% =========================================================

% ---------------------------------------------------------
% 1. Paket
% ---------------------------------------------------------
\usepackage[T1]{fontenc}
\usepackage[utf8]{inputenc}
\usepackage[indonesian]{babel}
\usepackage{lmodern}
\usepackage{tikz}
\usetikzlibrary{arrows.meta, positioning, calc, shapes.geometric}
\usepackage{booktabs}
\usepackage{array}
\usepackage{amsmath}
\usepackage{amsfonts}
\usepackage{pifont}          % \ding untuk tanda checklist
\usepackage{listings}

% ---------------------------------------------------------
% 2. Palet warna OBE (identik dengan buku)
% ---------------------------------------------------------
\definecolor{obeBlue}{HTML}{1F3864}
\definecolor{obeBlueLight}{HTML}{DCE4F5}
\definecolor{obeTeal}{HTML}{0F6E6E}
\definecolor{obeTealLight}{HTML}{D9F0EE}
\definecolor{obeGreen}{HTML}{2E6B34}
\definecolor{obeGreenLight}{HTML}{DEEFDF}
\definecolor{obeAmber}{HTML}{8A5A00}
\definecolor{obeAmberLight}{HTML}{FDF1D6}
\definecolor{obeMaroon}{HTML}{7A2233}
\definecolor{obeMaroonLight}{HTML}{F7E1E5}
\definecolor{obeGray}{HTML}{4A4A4A}
\definecolor{obeGrayLight}{HTML}{F0F0F0}
\definecolor{codeBack}{HTML}{F7F7F4}
\definecolor{codeFrame}{HTML}{C9C9C0}

% ---------------------------------------------------------
% 3. Identitas mata kuliah (dipakai seluruh deck)
% ---------------------------------------------------------
\newcommand{\NamaMK}{Bahasa Pemrograman Tingkat Rendah}
\newcommand{\NamaMKShort}{Bhs.\ Pemrograman Tingkat Rendah}
\newcommand{\KodeMK}{MWCA010}
\newcommand{\Prodi}{Program Studi Teknik Informatika}
\newcommand{\Fakultas}{Fakultas Teknologi Informasi}
\newcommand{\PerguruanTinggi}{Universitas Bale Bandung}
\newcommand{\DosenPengampu}{Cecep Suwanda, S.Si., M.Kom.}
\newcommand{\SemesterTeks}{Semester Ganjil}

% \PertemuanTeks dan \JudulBab ditetapkan di masing-masing
% berkas minggu-NN.tex sebelum \begin{document}.
\providecommand{\PertemuanTeks}{Pertemuan}

% ---------------------------------------------------------
% 4. Tema beamer
% ---------------------------------------------------------
\usetheme{default}
\setbeamertemplate{navigation symbols}{}

\setbeamercolor{normal text}{fg=black!88}
\setbeamercolor{structure}{fg=obeBlue}
\setbeamercolor{alerted text}{fg=obeMaroon}

% --- judul slide: pita biru penuh lebar ---
\setbeamercolor{frametitle}{fg=white,bg=obeBlue}
\setbeamerfont{frametitle}{series=\bfseries,size=\large}
\setbeamertemplate{frametitle}{%
    \nointerlineskip%
    \begin{beamercolorbox}[%
        wd=\paperwidth,ht=2.9ex,dp=1.5ex,
        leftskip=0.4cm,rightskip=0.4cm]{frametitle}%
        \usebeamerfont{frametitle}\insertframetitle%
    \end{beamercolorbox}%
    \nointerlineskip%
    \hbox{\textcolor{obeTeal}{\rule{\paperwidth}{1.1pt}}}%
}

% --- kaki slide: identitas mata kuliah, pertemuan, nomor ---
\setbeamercolor{footline}{fg=white,bg=obeGray}
\setbeamerfont{footline}{size=\scriptsize,family=\sffamily}
\setbeamertemplate{footline}{%
    \leavevmode%
    \hbox{%
        \begin{beamercolorbox}[wd=0.42\paperwidth,ht=2.1ex,dp=1.1ex,
            leftskip=0.4cm]{footline}%
            \usebeamerfont{footline}%
            \KodeMK\ \textbullet\ \NamaMKShort%
        \end{beamercolorbox}%
        \begin{beamercolorbox}[wd=0.32\paperwidth,ht=2.1ex,dp=1.1ex,
            center]{footline}%
            \usebeamerfont{footline}\PertemuanTeks%
        \end{beamercolorbox}%
        \begin{beamercolorbox}[wd=0.26\paperwidth,ht=2.1ex,dp=1.1ex,
            rightskip=0.4cm]{footline}%
            \usebeamerfont{footline}%
            \insertframenumber\,/\,\inserttotalframenumber%
        \end{beamercolorbox}%
    }%
    \vskip0pt%
}

% --- daftar ---
\setbeamercolor{itemize item}{fg=obeTeal}
\setbeamercolor{itemize subitem}{fg=obeTeal}
\setbeamercolor{enumerate item}{fg=obeBlue}
\setbeamertemplate{itemize item}{\raise0.2ex\hbox{\tiny$\blacksquare$}}
\setbeamertemplate{itemize subitem}{\raise0.2ex\hbox{\scriptsize$\bullet$}}

% --- kotak (pengganti tcolorbox pada buku) ---
\setbeamercolor{block title}{bg=obeBlueLight,fg=obeBlue}
\setbeamercolor{block body}{bg=obeBlueLight!45,fg=black!88}
\setbeamercolor{block title alerted}{bg=obeMaroonLight,fg=obeMaroon}
\setbeamercolor{block body alerted}{bg=obeMaroonLight!45,fg=black!88}
\setbeamercolor{block title example}{bg=obeTealLight,fg=obeTeal}
\setbeamercolor{block body example}{bg=obeTealLight!45,fg=black!88}

% ---------------------------------------------------------
% 5. Makro notasi (sama persis dengan buku)
%
% Catatan: \note TIDAK didefinisikan di sini karena beamer
% sudah memakai \note untuk catatan pembicara. Untuk catatan
% di slide dipakai \catatan.
% ---------------------------------------------------------
\newcommand{\code}[1]{\texttt{#1}}
\newcommand{\asm}[1]{\texttt{#1}}
\newcommand{\reg}[1]{\texttt{#1}}
\newcommand{\flag}[1]{\texttt{#1}}
\newcommand{\dir}[1]{\texttt{#1}}
\newcommand{\file}[1]{\texttt{#1}}
\newcommand{\eng}[1]{\textit{#1}}

\newcommand{\hex}[1]{\texttt{#1h}}
\newcommand{\bin}[1]{\texttt{#1b}}
\newcommand{\dec}[1]{\texttt{#1}}

\newcommand{\important}[1]{\textbf{Penting:} #1}
\newcommand{\catatan}[1]{\textbf{Catatan:} #1}
\newcommand{\term}[1]{\textbf{#1}}

\newcommand{\CPL}[1]{\textbf{CPL#1}}
\newcommand{\CPMK}[1]{\textbf{CPMK#1}}
\newcommand{\SubCPMK}[1]{\textbf{Sub-CPMK #1}}

\newcommand{\checkbox}{\ding{113}\hspace{0.35em}}
\newcommand{\checked}{\ding{51}\hspace{0.35em}}
\newcommand{\crossed}{\ding{55}\hspace{0.35em}}

\newcommand{\flowarrow}{\ensuremath{\rightarrow}}

% ---------------------------------------------------------
% 6. Makro tata letak slide
% ---------------------------------------------------------

% Keterangan sumber materi (ketertelusuran ke buku)
\newcommand{\Sumber}[1]{%
    \vfill
    {\scriptsize\color{obeGray}Sumber: #1\par}%
}

% Frame judul deck: \JudulDeck{<judul bab>}{<pertemuan>}
\newcommand{\JudulDeck}[2]{%
\begin{frame}[plain]
    \vspace{0.25cm}
    {\color{obeBlue}\sffamily\bfseries\small \NamaMK\ (\KodeMK)\par}
    \vspace{0.12cm}
    {\color{obeGray}\sffamily\scriptsize
     \Prodi\ \textbullet\ \Fakultas\ \textbullet\ \PerguruanTinggi\par}
    \vfill
    {\color{obeBlue}\sffamily\LARGE\bfseries #1\par}
    \vspace{0.4cm}
    {\color{obeTeal}\sffamily\large #2\par}
    \vspace{0.45cm}
    {\color{obeTeal}\rule{\textwidth}{1pt}}
    \vfill
    {\color{obeGray}\sffamily\small Dosen Pengampu\par}
    \vspace{0.08cm}
    {\sffamily\bfseries \DosenPengampu\par}
    \vspace{0.2cm}
    {\color{obeGray}\sffamily\small \SemesterTeks\par}
    \vspace{0.25cm}
\end{frame}
}

% Frame pemisah bagian
\newcommand{\PemisahBagian}[1]{%
\begin{frame}[plain]
    \vfill
    \begin{center}
        {\color{obeTeal}\sffamily\small\bfseries BAGIAN\par}
        \vspace{0.35cm}
        {\color{obeBlue}\sffamily\Large\bfseries #1\par}
        \vspace{0.35cm}
        {\color{obeTeal}\rule{0.55\textwidth}{0.9pt}}
    \end{center}
    \vfill
\end{frame}
}

% ---------------------------------------------------------
% 7. Gaya penampil kode
% ---------------------------------------------------------
\input{listing-ppt}

%
% Palet warna disalin dari buku_ajar/config/styles.tex agar
% identitas visual slide sama dengan buku ajar.
% =========================================================

% ---------------------------------------------------------
% 1. Paket
% ---------------------------------------------------------
\usepackage[T1]{fontenc}
\usepackage[utf8]{inputenc}
\usepackage[indonesian]{babel}
\usepackage{lmodern}
\usepackage{tikz}
\usetikzlibrary{arrows.meta, positioning, calc, shapes.geometric}
\usepackage{booktabs}
\usepackage{array}
\usepackage{amsmath}
\usepackage{amsfonts}
\usepackage{pifont}          % \ding untuk tanda checklist
\usepackage{listings}

% ---------------------------------------------------------
% 2. Palet warna OBE (identik dengan buku)
% ---------------------------------------------------------
\definecolor{obeBlue}{HTML}{1F3864}
\definecolor{obeBlueLight}{HTML}{DCE4F5}
\definecolor{obeTeal}{HTML}{0F6E6E}
\definecolor{obeTealLight}{HTML}{D9F0EE}
\definecolor{obeGreen}{HTML}{2E6B34}
\definecolor{obeGreenLight}{HTML}{DEEFDF}
\definecolor{obeAmber}{HTML}{8A5A00}
\definecolor{obeAmberLight}{HTML}{FDF1D6}
\definecolor{obeMaroon}{HTML}{7A2233}
\definecolor{obeMaroonLight}{HTML}{F7E1E5}
\definecolor{obeGray}{HTML}{4A4A4A}
\definecolor{obeGrayLight}{HTML}{F0F0F0}
\definecolor{codeBack}{HTML}{F7F7F4}
\definecolor{codeFrame}{HTML}{C9C9C0}

% ---------------------------------------------------------
% 3. Identitas mata kuliah (dipakai seluruh deck)
% ---------------------------------------------------------
\newcommand{\NamaMK}{Bahasa Pemrograman Tingkat Rendah}
\newcommand{\NamaMKShort}{Bhs.\ Pemrograman Tingkat Rendah}
\newcommand{\KodeMK}{MWCA010}
\newcommand{\Prodi}{Program Studi Teknik Informatika}
\newcommand{\Fakultas}{Fakultas Teknologi Informasi}
\newcommand{\PerguruanTinggi}{Universitas Bale Bandung}
\newcommand{\DosenPengampu}{Cecep Suwanda, S.Si., M.Kom.}
\newcommand{\SemesterTeks}{Semester Ganjil}

% \PertemuanTeks dan \JudulBab ditetapkan di masing-masing
% berkas minggu-NN.tex sebelum \begin{document}.
\providecommand{\PertemuanTeks}{Pertemuan}

% ---------------------------------------------------------
% 4. Tema beamer
% ---------------------------------------------------------
\usetheme{default}
\setbeamertemplate{navigation symbols}{}

\setbeamercolor{normal text}{fg=black!88}
\setbeamercolor{structure}{fg=obeBlue}
\setbeamercolor{alerted text}{fg=obeMaroon}

% --- judul slide: pita biru penuh lebar ---
\setbeamercolor{frametitle}{fg=white,bg=obeBlue}
\setbeamerfont{frametitle}{series=\bfseries,size=\large}
\setbeamertemplate{frametitle}{%
    \nointerlineskip%
    \begin{beamercolorbox}[%
        wd=\paperwidth,ht=2.9ex,dp=1.5ex,
        leftskip=0.4cm,rightskip=0.4cm]{frametitle}%
        \usebeamerfont{frametitle}\insertframetitle%
    \end{beamercolorbox}%
    \nointerlineskip%
    \hbox{\textcolor{obeTeal}{\rule{\paperwidth}{1.1pt}}}%
}

% --- kaki slide: identitas mata kuliah, pertemuan, nomor ---
\setbeamercolor{footline}{fg=white,bg=obeGray}
\setbeamerfont{footline}{size=\scriptsize,family=\sffamily}
\setbeamertemplate{footline}{%
    \leavevmode%
    \hbox{%
        \begin{beamercolorbox}[wd=0.42\paperwidth,ht=2.1ex,dp=1.1ex,
            leftskip=0.4cm]{footline}%
            \usebeamerfont{footline}%
            \KodeMK\ \textbullet\ \NamaMKShort%
        \end{beamercolorbox}%
        \begin{beamercolorbox}[wd=0.32\paperwidth,ht=2.1ex,dp=1.1ex,
            center]{footline}%
            \usebeamerfont{footline}\PertemuanTeks%
        \end{beamercolorbox}%
        \begin{beamercolorbox}[wd=0.26\paperwidth,ht=2.1ex,dp=1.1ex,
            rightskip=0.4cm]{footline}%
            \usebeamerfont{footline}%
            \insertframenumber\,/\,\inserttotalframenumber%
        \end{beamercolorbox}%
    }%
    \vskip0pt%
}

% --- daftar ---
\setbeamercolor{itemize item}{fg=obeTeal}
\setbeamercolor{itemize subitem}{fg=obeTeal}
\setbeamercolor{enumerate item}{fg=obeBlue}
\setbeamertemplate{itemize item}{\raise0.2ex\hbox{\tiny$\blacksquare$}}
\setbeamertemplate{itemize subitem}{\raise0.2ex\hbox{\scriptsize$\bullet$}}

% --- kotak (pengganti tcolorbox pada buku) ---
\setbeamercolor{block title}{bg=obeBlueLight,fg=obeBlue}
\setbeamercolor{block body}{bg=obeBlueLight!45,fg=black!88}
\setbeamercolor{block title alerted}{bg=obeMaroonLight,fg=obeMaroon}
\setbeamercolor{block body alerted}{bg=obeMaroonLight!45,fg=black!88}
\setbeamercolor{block title example}{bg=obeTealLight,fg=obeTeal}
\setbeamercolor{block body example}{bg=obeTealLight!45,fg=black!88}

% ---------------------------------------------------------
% 5. Makro notasi (sama persis dengan buku)
%
% Catatan: \note TIDAK didefinisikan di sini karena beamer
% sudah memakai \note untuk catatan pembicara. Untuk catatan
% di slide dipakai \catatan.
% ---------------------------------------------------------
\newcommand{\code}[1]{\texttt{#1}}
\newcommand{\asm}[1]{\texttt{#1}}
\newcommand{\reg}[1]{\texttt{#1}}
\newcommand{\flag}[1]{\texttt{#1}}
\newcommand{\dir}[1]{\texttt{#1}}
\newcommand{\file}[1]{\texttt{#1}}
\newcommand{\eng}[1]{\textit{#1}}

\newcommand{\hex}[1]{\texttt{#1h}}
\newcommand{\bin}[1]{\texttt{#1b}}
\newcommand{\dec}[1]{\texttt{#1}}

\newcommand{\important}[1]{\textbf{Penting:} #1}
\newcommand{\catatan}[1]{\textbf{Catatan:} #1}
\newcommand{\term}[1]{\textbf{#1}}

\newcommand{\CPL}[1]{\textbf{CPL#1}}
\newcommand{\CPMK}[1]{\textbf{CPMK#1}}
\newcommand{\SubCPMK}[1]{\textbf{Sub-CPMK #1}}

\newcommand{\checkbox}{\ding{113}\hspace{0.35em}}
\newcommand{\checked}{\ding{51}\hspace{0.35em}}
\newcommand{\crossed}{\ding{55}\hspace{0.35em}}

\newcommand{\flowarrow}{\ensuremath{\rightarrow}}

% ---------------------------------------------------------
% 6. Makro tata letak slide
% ---------------------------------------------------------

% Keterangan sumber materi (ketertelusuran ke buku)
\newcommand{\Sumber}[1]{%
    \vfill
    {\scriptsize\color{obeGray}Sumber: #1\par}%
}

% Frame judul deck: \JudulDeck{<judul bab>}{<pertemuan>}
\newcommand{\JudulDeck}[2]{%
\begin{frame}[plain]
    \vspace{0.25cm}
    {\color{obeBlue}\sffamily\bfseries\small \NamaMK\ (\KodeMK)\par}
    \vspace{0.12cm}
    {\color{obeGray}\sffamily\scriptsize
     \Prodi\ \textbullet\ \Fakultas\ \textbullet\ \PerguruanTinggi\par}
    \vfill
    {\color{obeBlue}\sffamily\LARGE\bfseries #1\par}
    \vspace{0.4cm}
    {\color{obeTeal}\sffamily\large #2\par}
    \vspace{0.45cm}
    {\color{obeTeal}\rule{\textwidth}{1pt}}
    \vfill
    {\color{obeGray}\sffamily\small Dosen Pengampu\par}
    \vspace{0.08cm}
    {\sffamily\bfseries \DosenPengampu\par}
    \vspace{0.2cm}
    {\color{obeGray}\sffamily\small \SemesterTeks\par}
    \vspace{0.25cm}
\end{frame}
}

% Frame pemisah bagian
\newcommand{\PemisahBagian}[1]{%
\begin{frame}[plain]
    \vfill
    \begin{center}
        {\color{obeTeal}\sffamily\small\bfseries BAGIAN\par}
        \vspace{0.35cm}
        {\color{obeBlue}\sffamily\Large\bfseries #1\par}
        \vspace{0.35cm}
        {\color{obeTeal}\rule{0.55\textwidth}{0.9pt}}
    \end{center}
    \vfill
\end{frame}
}

% ---------------------------------------------------------
% 7. Gaya penampil kode
% ---------------------------------------------------------
% =========================================================
% ppt/src/listing-ppt.tex
% Gaya penampil kode untuk slide.
%
% Diturunkan dari buku_ajar/config/listings.tex, tetapi
% disesuaikan untuk slide: ukuran lebih kecil (karena frame
% 16:9 dan kode harus utuh terlihat) dan jarak lebih rapat.
% =========================================================

% ---------------------------------------------------------
% 1. Definisi bahasa MASM/TASM 8086 (salinan dari buku)
% ---------------------------------------------------------
\lstdefinelanguage{MASM}{
    morekeywords={
        MOV, XCHG, LEA, LDS, LES, PUSH, POP, PUSHA, POPA, PUSHF, POPF,
        LAHF, SAHF, XLAT, IN, OUT, CBW, CWD,
        ADD, ADC, SUB, SBB, INC, DEC, NEG, CMP, MUL, IMUL, DIV, IDIV,
        AAA, AAS, AAM, AAD, DAA, DAS,
        AND, OR, XOR, NOT, TEST, SHL, SHR, SAL, SAR, ROL, ROR, RCL, RCR,
        JMP, JE, JZ, JNE, JNZ, JG, JGE, JL, JLE, JA, JAE, JB, JBE,
        JC, JNC, JS, JNS, JO, JNO, JP, JPE, JNP, JPO, JCXZ,
        LOOP, LOOPE, LOOPZ, LOOPNE, LOOPNZ, CALL, RET, INT, IRET, INTO,
        HLT, NOP, CLI, STI, CLD, STD, CLC, STC, CMC, WAIT, ESC, LOCK,
        DB, DW, DD, DQ, DT, EQU, ORG, SEGMENT, ENDS, ASSUME, END,
        PROC, ENDP, MACRO, ENDM, LOCAL, PUBLIC, EXTRN, EXTERN,
        INCLUDE, INCLUDELIB, MODEL, STACK, CODE, DATA, CONST, ALIGN, EVEN,
        OFFSET, PTR, BYTE, WORD, DWORD, QWORD, TBYTE,
        NEAR, FAR, SHORT, LABEL, TYPE, SIZE, LENGTH, DUP, THIS,
        AX, BX, CX, DX, AL, AH, BL, BH, CL, CH, DL, DH,
        SP, BP, SI, DI, IP, CS, DS, SS, ES, FLAGS,
        CF, PF, AF, ZF, TF, IF, DF, OF
    },
    sensitive=false,
    morecomment=[l]{;},
    morestring=[b]',
    morestring=[b]"
}

% ---------------------------------------------------------
% 2. Gaya kode assembly untuk slide
% ---------------------------------------------------------
\lstdefinestyle{asmslide}{
    language=MASM,
    basicstyle=\ttfamily\scriptsize,
    keywordstyle=\color{obeBlue}\bfseries,
    commentstyle=\color{obeGreen}\itshape,
    stringstyle=\color{obeMaroon},
    numberstyle=\tiny\color{obeGray},
    numbers=left,
    numbersep=5pt,
    stepnumber=1,
    backgroundcolor=\color{codeBack},
    frame=single,
    rulecolor=\color{codeFrame},
    framesep=4pt,
    breaklines=true,
    breakatwhitespace=false,
    % CATATAN PENTING: opsi `postbreak` TIDAK boleh dipakai di sini.
    % Di beamer, `breaklines=true` bersama `postbreak=<kode>` membuat
    % listings gagal dirakit dengan galat
    %   "Missing number, treated as zero <to be read again> \kern".
    % Karena \lstset{style=asmslide} bersifat global, opsi postbreak
    % juga bocor ke gaya lain (mis. outslide/keluaran) dan merusaknya.
    % Di buku ajar (kelas `book`) justru sebaliknya: postbreak aman dan
    % dipakai. Jadi larangan ini khusus untuk slide.
    showstringspaces=false,
    tabsize=4,
    columns=fullflexible,
    keepspaces=true,
    upquote=true,
    xleftmargin=12pt,
    aboveskip=0.4em,
    belowskip=0.4em
}

% ---------------------------------------------------------
% 3. Gaya keluaran program / sesi debugger untuk slide
% ---------------------------------------------------------
\lstdefinestyle{outslide}{
    basicstyle=\ttfamily\scriptsize,
    keywordstyle=\color{black},
    commentstyle=\color{obeGray},
    stringstyle=\color{black},
    numbers=none,
    backgroundcolor=\color{obeGrayLight},
    frame=single,
    rulecolor=\color{codeFrame},
    framesep=4pt,
    breaklines=true,
    showstringspaces=false,
    columns=fullflexible,
    keepspaces=true,
    upquote=true,
    xleftmargin=4pt,
    aboveskip=0.4em,
    belowskip=0.4em
}

\lstset{style=asmslide}

% ---------------------------------------------------------
% 4. Environment ringkas untuk dipakai di dalam frame
%
% Frame yang memuat kedua environment ini WAJIB memakai
% opsi [fragile], karena isinya verbatim.
% ---------------------------------------------------------
\lstnewenvironment{kodeasm}[1][]{%
    \lstset{style=asmslide, #1}%
}{}

\lstnewenvironment{keluaran}[1][]{%
    \lstset{style=outslide, #1}%
}{}


%
% Palet warna disalin dari buku_ajar/config/styles.tex agar
% identitas visual slide sama dengan buku ajar.
% =========================================================

% ---------------------------------------------------------
% 1. Paket
% ---------------------------------------------------------
\usepackage[T1]{fontenc}
\usepackage[utf8]{inputenc}
\usepackage[indonesian]{babel}
\usepackage{lmodern}
\usepackage{tikz}
\usetikzlibrary{arrows.meta, positioning, calc, shapes.geometric}
\usepackage{booktabs}
\usepackage{array}
\usepackage{amsmath}
\usepackage{amsfonts}
\usepackage{pifont}          % \ding untuk tanda checklist
\usepackage{listings}

% ---------------------------------------------------------
% 2. Palet warna OBE (identik dengan buku)
% ---------------------------------------------------------
\definecolor{obeBlue}{HTML}{1F3864}
\definecolor{obeBlueLight}{HTML}{DCE4F5}
\definecolor{obeTeal}{HTML}{0F6E6E}
\definecolor{obeTealLight}{HTML}{D9F0EE}
\definecolor{obeGreen}{HTML}{2E6B34}
\definecolor{obeGreenLight}{HTML}{DEEFDF}
\definecolor{obeAmber}{HTML}{8A5A00}
\definecolor{obeAmberLight}{HTML}{FDF1D6}
\definecolor{obeMaroon}{HTML}{7A2233}
\definecolor{obeMaroonLight}{HTML}{F7E1E5}
\definecolor{obeGray}{HTML}{4A4A4A}
\definecolor{obeGrayLight}{HTML}{F0F0F0}
\definecolor{codeBack}{HTML}{F7F7F4}
\definecolor{codeFrame}{HTML}{C9C9C0}

% ---------------------------------------------------------
% 3. Identitas mata kuliah (dipakai seluruh deck)
% ---------------------------------------------------------
\newcommand{\NamaMK}{Bahasa Pemrograman Tingkat Rendah}
\newcommand{\NamaMKShort}{Bhs.\ Pemrograman Tingkat Rendah}
\newcommand{\KodeMK}{MWCA010}
\newcommand{\Prodi}{Program Studi Teknik Informatika}
\newcommand{\Fakultas}{Fakultas Teknologi Informasi}
\newcommand{\PerguruanTinggi}{Universitas Bale Bandung}
\newcommand{\DosenPengampu}{Cecep Suwanda, S.Si., M.Kom.}
\newcommand{\SemesterTeks}{Semester Ganjil}

% \PertemuanTeks dan \JudulBab ditetapkan di masing-masing
% berkas minggu-NN.tex sebelum \begin{document}.
\providecommand{\PertemuanTeks}{Pertemuan}

% ---------------------------------------------------------
% 4. Tema beamer
% ---------------------------------------------------------
\usetheme{default}
\setbeamertemplate{navigation symbols}{}

\setbeamercolor{normal text}{fg=black!88}
\setbeamercolor{structure}{fg=obeBlue}
\setbeamercolor{alerted text}{fg=obeMaroon}

% --- judul slide: pita biru penuh lebar ---
\setbeamercolor{frametitle}{fg=white,bg=obeBlue}
\setbeamerfont{frametitle}{series=\bfseries,size=\large}
\setbeamertemplate{frametitle}{%
    \nointerlineskip%
    \begin{beamercolorbox}[%
        wd=\paperwidth,ht=2.9ex,dp=1.5ex,
        leftskip=0.4cm,rightskip=0.4cm]{frametitle}%
        \usebeamerfont{frametitle}\insertframetitle%
    \end{beamercolorbox}%
    \nointerlineskip%
    \hbox{\textcolor{obeTeal}{\rule{\paperwidth}{1.1pt}}}%
}

% --- kaki slide: identitas mata kuliah, pertemuan, nomor ---
\setbeamercolor{footline}{fg=white,bg=obeGray}
\setbeamerfont{footline}{size=\scriptsize,family=\sffamily}
\setbeamertemplate{footline}{%
    \leavevmode%
    \hbox{%
        \begin{beamercolorbox}[wd=0.42\paperwidth,ht=2.1ex,dp=1.1ex,
            leftskip=0.4cm]{footline}%
            \usebeamerfont{footline}%
            \KodeMK\ \textbullet\ \NamaMKShort%
        \end{beamercolorbox}%
        \begin{beamercolorbox}[wd=0.32\paperwidth,ht=2.1ex,dp=1.1ex,
            center]{footline}%
            \usebeamerfont{footline}\PertemuanTeks%
        \end{beamercolorbox}%
        \begin{beamercolorbox}[wd=0.26\paperwidth,ht=2.1ex,dp=1.1ex,
            rightskip=0.4cm]{footline}%
            \usebeamerfont{footline}%
            \insertframenumber\,/\,\inserttotalframenumber%
        \end{beamercolorbox}%
    }%
    \vskip0pt%
}

% --- daftar ---
\setbeamercolor{itemize item}{fg=obeTeal}
\setbeamercolor{itemize subitem}{fg=obeTeal}
\setbeamercolor{enumerate item}{fg=obeBlue}
\setbeamertemplate{itemize item}{\raise0.2ex\hbox{\tiny$\blacksquare$}}
\setbeamertemplate{itemize subitem}{\raise0.2ex\hbox{\scriptsize$\bullet$}}

% --- kotak (pengganti tcolorbox pada buku) ---
\setbeamercolor{block title}{bg=obeBlueLight,fg=obeBlue}
\setbeamercolor{block body}{bg=obeBlueLight!45,fg=black!88}
\setbeamercolor{block title alerted}{bg=obeMaroonLight,fg=obeMaroon}
\setbeamercolor{block body alerted}{bg=obeMaroonLight!45,fg=black!88}
\setbeamercolor{block title example}{bg=obeTealLight,fg=obeTeal}
\setbeamercolor{block body example}{bg=obeTealLight!45,fg=black!88}

% ---------------------------------------------------------
% 5. Makro notasi (sama persis dengan buku)
%
% Catatan: \note TIDAK didefinisikan di sini karena beamer
% sudah memakai \note untuk catatan pembicara. Untuk catatan
% di slide dipakai \catatan.
% ---------------------------------------------------------
\newcommand{\code}[1]{\texttt{#1}}
\newcommand{\asm}[1]{\texttt{#1}}
\newcommand{\reg}[1]{\texttt{#1}}
\newcommand{\flag}[1]{\texttt{#1}}
\newcommand{\dir}[1]{\texttt{#1}}
\newcommand{\file}[1]{\texttt{#1}}
\newcommand{\eng}[1]{\textit{#1}}

\newcommand{\hex}[1]{\texttt{#1h}}
\newcommand{\bin}[1]{\texttt{#1b}}
\newcommand{\dec}[1]{\texttt{#1}}

\newcommand{\important}[1]{\textbf{Penting:} #1}
\newcommand{\catatan}[1]{\textbf{Catatan:} #1}
\newcommand{\term}[1]{\textbf{#1}}

\newcommand{\CPL}[1]{\textbf{CPL#1}}
\newcommand{\CPMK}[1]{\textbf{CPMK#1}}
\newcommand{\SubCPMK}[1]{\textbf{Sub-CPMK #1}}

\newcommand{\checkbox}{\ding{113}\hspace{0.35em}}
\newcommand{\checked}{\ding{51}\hspace{0.35em}}
\newcommand{\crossed}{\ding{55}\hspace{0.35em}}

\newcommand{\flowarrow}{\ensuremath{\rightarrow}}

% ---------------------------------------------------------
% 6. Makro tata letak slide
% ---------------------------------------------------------

% Keterangan sumber materi (ketertelusuran ke buku)
\newcommand{\Sumber}[1]{%
    \vfill
    {\scriptsize\color{obeGray}Sumber: #1\par}%
}

% Frame judul deck: \JudulDeck{<judul bab>}{<pertemuan>}
\newcommand{\JudulDeck}[2]{%
\begin{frame}[plain]
    \vspace{0.25cm}
    {\color{obeBlue}\sffamily\bfseries\small \NamaMK\ (\KodeMK)\par}
    \vspace{0.12cm}
    {\color{obeGray}\sffamily\scriptsize
     \Prodi\ \textbullet\ \Fakultas\ \textbullet\ \PerguruanTinggi\par}
    \vfill
    {\color{obeBlue}\sffamily\LARGE\bfseries #1\par}
    \vspace{0.4cm}
    {\color{obeTeal}\sffamily\large #2\par}
    \vspace{0.45cm}
    {\color{obeTeal}\rule{\textwidth}{1pt}}
    \vfill
    {\color{obeGray}\sffamily\small Dosen Pengampu\par}
    \vspace{0.08cm}
    {\sffamily\bfseries \DosenPengampu\par}
    \vspace{0.2cm}
    {\color{obeGray}\sffamily\small \SemesterTeks\par}
    \vspace{0.25cm}
\end{frame}
}

% Frame pemisah bagian
\newcommand{\PemisahBagian}[1]{%
\begin{frame}[plain]
    \vfill
    \begin{center}
        {\color{obeTeal}\sffamily\small\bfseries BAGIAN\par}
        \vspace{0.35cm}
        {\color{obeBlue}\sffamily\Large\bfseries #1\par}
        \vspace{0.35cm}
        {\color{obeTeal}\rule{0.55\textwidth}{0.9pt}}
    \end{center}
    \vfill
\end{frame}
}

% ---------------------------------------------------------
% 7. Gaya penampil kode
% ---------------------------------------------------------
% =========================================================
% ppt/src/listing-ppt.tex
% Gaya penampil kode untuk slide.
%
% Diturunkan dari buku_ajar/config/listings.tex, tetapi
% disesuaikan untuk slide: ukuran lebih kecil (karena frame
% 16:9 dan kode harus utuh terlihat) dan jarak lebih rapat.
% =========================================================

% ---------------------------------------------------------
% 1. Definisi bahasa MASM/TASM 8086 (salinan dari buku)
% ---------------------------------------------------------
\lstdefinelanguage{MASM}{
    morekeywords={
        MOV, XCHG, LEA, LDS, LES, PUSH, POP, PUSHA, POPA, PUSHF, POPF,
        LAHF, SAHF, XLAT, IN, OUT, CBW, CWD,
        ADD, ADC, SUB, SBB, INC, DEC, NEG, CMP, MUL, IMUL, DIV, IDIV,
        AAA, AAS, AAM, AAD, DAA, DAS,
        AND, OR, XOR, NOT, TEST, SHL, SHR, SAL, SAR, ROL, ROR, RCL, RCR,
        JMP, JE, JZ, JNE, JNZ, JG, JGE, JL, JLE, JA, JAE, JB, JBE,
        JC, JNC, JS, JNS, JO, JNO, JP, JPE, JNP, JPO, JCXZ,
        LOOP, LOOPE, LOOPZ, LOOPNE, LOOPNZ, CALL, RET, INT, IRET, INTO,
        HLT, NOP, CLI, STI, CLD, STD, CLC, STC, CMC, WAIT, ESC, LOCK,
        DB, DW, DD, DQ, DT, EQU, ORG, SEGMENT, ENDS, ASSUME, END,
        PROC, ENDP, MACRO, ENDM, LOCAL, PUBLIC, EXTRN, EXTERN,
        INCLUDE, INCLUDELIB, MODEL, STACK, CODE, DATA, CONST, ALIGN, EVEN,
        OFFSET, PTR, BYTE, WORD, DWORD, QWORD, TBYTE,
        NEAR, FAR, SHORT, LABEL, TYPE, SIZE, LENGTH, DUP, THIS,
        AX, BX, CX, DX, AL, AH, BL, BH, CL, CH, DL, DH,
        SP, BP, SI, DI, IP, CS, DS, SS, ES, FLAGS,
        CF, PF, AF, ZF, TF, IF, DF, OF
    },
    sensitive=false,
    morecomment=[l]{;},
    morestring=[b]',
    morestring=[b]"
}

% ---------------------------------------------------------
% 2. Gaya kode assembly untuk slide
% ---------------------------------------------------------
\lstdefinestyle{asmslide}{
    language=MASM,
    basicstyle=\ttfamily\scriptsize,
    keywordstyle=\color{obeBlue}\bfseries,
    commentstyle=\color{obeGreen}\itshape,
    stringstyle=\color{obeMaroon},
    numberstyle=\tiny\color{obeGray},
    numbers=left,
    numbersep=5pt,
    stepnumber=1,
    backgroundcolor=\color{codeBack},
    frame=single,
    rulecolor=\color{codeFrame},
    framesep=4pt,
    breaklines=true,
    breakatwhitespace=false,
    % CATATAN PENTING: opsi `postbreak` TIDAK boleh dipakai di sini.
    % Di beamer, `breaklines=true` bersama `postbreak=<kode>` membuat
    % listings gagal dirakit dengan galat
    %   "Missing number, treated as zero <to be read again> \kern".
    % Karena \lstset{style=asmslide} bersifat global, opsi postbreak
    % juga bocor ke gaya lain (mis. outslide/keluaran) dan merusaknya.
    % Di buku ajar (kelas `book`) justru sebaliknya: postbreak aman dan
    % dipakai. Jadi larangan ini khusus untuk slide.
    showstringspaces=false,
    tabsize=4,
    columns=fullflexible,
    keepspaces=true,
    upquote=true,
    xleftmargin=12pt,
    aboveskip=0.4em,
    belowskip=0.4em
}

% ---------------------------------------------------------
% 3. Gaya keluaran program / sesi debugger untuk slide
% ---------------------------------------------------------
\lstdefinestyle{outslide}{
    basicstyle=\ttfamily\scriptsize,
    keywordstyle=\color{black},
    commentstyle=\color{obeGray},
    stringstyle=\color{black},
    numbers=none,
    backgroundcolor=\color{obeGrayLight},
    frame=single,
    rulecolor=\color{codeFrame},
    framesep=4pt,
    breaklines=true,
    showstringspaces=false,
    columns=fullflexible,
    keepspaces=true,
    upquote=true,
    xleftmargin=4pt,
    aboveskip=0.4em,
    belowskip=0.4em
}

\lstset{style=asmslide}

% ---------------------------------------------------------
% 4. Environment ringkas untuk dipakai di dalam frame
%
% Frame yang memuat kedua environment ini WAJIB memakai
% opsi [fragile], karena isinya verbatim.
% ---------------------------------------------------------
\lstnewenvironment{kodeasm}[1][]{%
    \lstset{style=asmslide, #1}%
}{}

\lstnewenvironment{keluaran}[1][]{%
    \lstset{style=outslide, #1}%
}{}



\renewcommand{\PertemuanTeks}{Pertemuan 15}

\begin{document}

% =========================================================
% 1. JUDUL
% =========================================================
\JudulDeck{Debugging Level Instruksi}{Pertemuan 15}

% =========================================================
% 2. CAPAIAN DAN TUJUAN
% =========================================================
\begin{frame}{Capaian dan Tujuan Pembelajaran}
\begin{block}{Sub-CPMK 4.1--4.3}
\footnotesize
\textbf{4.1} Melakukan penelusuran program per instruksi
(\eng{single step}/\eng{trace}) dan memantau perubahan register
serta \eng{flag}.\\
\textbf{4.2} Memeriksa isi memori dan stack (\eng{dump}) untuk
menemukan kesalahan sintaks dan kesalahan logika.\\
\textbf{4.3} Menganalisis penyebab program berhenti (\eng{hang}) dan
memperbaiki kesalahan berdasarkan hasil penelusuran.
\end{block}

\vspace{0.1em}
{\small\textbf{CPMK yang didukung:} \CPMK{04}\quad\textbullet\quad
\textbf{CPL yang didukung:} \CPL{03}\par}

\vspace{0.15em}
\begin{exampleblock}{Setelah pertemuan ini, mahasiswa mampu:}
\footnotesize
\begin{enumerate}
    \item membedakan kesalahan sintaks, logika, dan \eng{hang};
    \item menggunakan DEBUG, emu8086, dan Turbo Debugger;
    \item mencatat register dan \eng{flag} per instruksi;
    \item memeriksa memori dan stack dengan \eng{dump};
    \item memperbaiki perulangan tak berujung dan stack tidak
          seimbang.
\end{enumerate}
\end{exampleblock}
\end{frame}

% =========================================================
% 3. PETA MATERI
% =========================================================
\begin{frame}{Peta Materi}
\small
\begin{enumerate}
    \item \textbf{Jenis Kesalahan dan Alat Debugging}\\
          {\footnotesize sintaks, logika, hang; DEBUG, emu8086,
          Turbo Debugger; alur debugging}

    \vspace{0.25em}
    \item \textbf{Menelusuri Program per Instruksi}\\
          {\footnotesize menyiapkan trace; register dan flag;
          perulangan tak berujung}

    \vspace{0.25em}
    \item \textbf{Memeriksa Memori, Stack, dan Memperbaiki Hang}\\
          {\footnotesize dump; stack tidak seimbang; stack frame;
          analisis hang}
\end{enumerate}

\vspace{0.35em}
\begin{block}{Praktik pada pertemuan ini}
\small Menemukan dan memperbaiki kesalahan dengan \eng{single step}
dan \eng{dump}.
\end{block}
\end{frame}

% =========================================================
% BAGIAN 1
% =========================================================
\PemisahBagian{Jenis Kesalahan dan Alat Debugging}

% ---------------------------------------------------------
\begin{frame}{Tiga Jenis Kesalahan}
\footnotesize
\begin{center}
\begin{tabular}{@{}>{\raggedright\arraybackslash}p{2.8cm}
    >{\raggedright\arraybackslash}p{5.0cm}
    >{\raggedright\arraybackslash}p{5.0cm}@{}}
    \toprule
    \textbf{Jenis} & \textbf{Gejala} & \textbf{Alat utama} \\
    \midrule
    Sintaks & Pesan TASM; tidak ada \file{.OBJ} &
    Daftar kesalahan assembler \\
    \addlinespace
    Logika & Program selesai, hasil salah &
    Single step, pemeriksaan register \\
    \addlinespace
    \eng{Hang} & Program tidak kembali ke DOS &
    Single step, pemeriksaan stack \\
    \bottomrule
\end{tabular}
\end{center}
\Sumber{Bab XVI, Bagian 1}
\end{frame}

% ---------------------------------------------------------
\begin{frame}{Alat Debugging}
\small
\begin{itemize}
    \item \textbf{DEBUG}: \eng{debugger} baris perintah DOS dengan
          perintah R, T, P, D, dan G.

    \vspace{0.18em}
    \item \textbf{emu8086}: tombol \eng{Single Step} serta jendela
          register, \eng{flag}, dan memori.

    \vspace{0.18em}
    \item \textbf{Turbo Debugger}: tampilan register, \eng{flag},
          memori, dan stack; F7 masuk instruksi, F8 melewati
          prosedur.
\end{itemize}

\vspace{0.2em}
\begin{block}{Prinsip}
\small Alat dipilih sesuai jenis kesalahan; hasil pengamatan dicatat,
bukan ditebak.
\end{block}
\Sumber{Bab XVI, Bagian 1}
\end{frame}

% ---------------------------------------------------------
\begin{frame}{Perintah DEBUG}
\footnotesize
\begin{center}
\begin{tabular}{@{}>{\raggedright\arraybackslash}p{1.8cm}
    >{\raggedright\arraybackslash}p{5.6cm}
    >{\raggedright\arraybackslash}p{5.4cm}@{}}
    \toprule
    \textbf{Perintah} & \textbf{Fungsi} & \textbf{Pemakaian} \\
    \midrule
    R & Menampilkan register dan flag &
    Keadaan sebelum/sesudah instruksi \\
    \addlinespace
    T & Menjalankan satu instruksi & Masuk ke prosedur \\
    \addlinespace
    P & Satu instruksi, melewati prosedur &
    Prosedur yang sudah dipercaya \\
    \addlinespace
    D & Menampilkan isi memori &
    Variabel, string, dan stack \\
    \addlinespace
    G & Menjalankan sampai alamat tertentu &
    Melewati bagian yang tidak dicurigai \\
    \bottomrule
\end{tabular}
\end{center}
\Sumber{Bab XVI, Bagian 1}
\end{frame}

% ---------------------------------------------------------
\begin{frame}{Alur Debugging}
\begin{center}
\begin{tikzpicture}[
    tahap/.style={
        draw=obeBlue, line width=0.8pt, fill=obeBlueLight,
        rounded corners=3pt, minimum width=2.6cm,
        minimum height=0.9cm, align=center, font=\small
    },
    panah/.style={-{Stealth[length=2.2mm]}, line width=0.9pt, obeTeal}
]
    \node[tahap] (rep) at (0,2.0) {Reproduksi};
    \node[tahap] (iso) at (3.0,2.0) {Isolasi};
    \node[tahap] (hip) at (6.0,2.0) {Hipotesis};
    \node[tahap] (uji) at (9.0,2.0) {Uji trace};
    \node[tahap] (fix) at (6.0,0.6) {Perbaiki};
    \node[tahap] (ver) at (3.0,0.6) {Verifikasi};

    \draw[panah] (rep) -- (iso);
    \draw[panah] (iso) -- (hip);
    \draw[panah] (hip) -- (uji);
    \draw[panah] (uji) -- (fix);
    \draw[panah] (fix) -- (ver);
\end{tikzpicture}
\end{center}

\vspace{0.2em}
{\footnotesize\color{obeGray}Perbaikan tanpa verifikasi belum
dianggap selesai.\par}
\Sumber{Bab XVI, Bagian 1}
\end{frame}

% ---------------------------------------------------------
\begin{frame}{Rangkuman Bagian 1}
\small
\begin{itemize}
    \item Tiga jenis kesalahan: sintaks, logika, dan \eng{hang}.

    \vspace{0.15em}
    \item Sintaks dilaporkan assembler; logika dan hang dicari dengan
          \eng{debugger}.

    \vspace{0.15em}
    \item DEBUG menyediakan R, T, P, D, dan G.

    \vspace{0.15em}
    \item emu8086 dan Turbo Debugger memberi tampilan register,
          flag, memori, dan stack.

    \vspace{0.15em}
    \item Debugging sistematis: reproduksi, isolasi, hipotesis, uji,
          perbaikan, verifikasi.
\end{itemize}
\end{frame}

% =========================================================
% BAGIAN 2
% =========================================================
\PemisahBagian{Menelusuri Program per Instruksi}

% ---------------------------------------------------------
\begin{frame}{Menyiapkan Penelusuran}
\small
\begin{enumerate}
    \item Rakit dan link program.

    \vspace{0.15em}
    \item Jalankan di bawah DEBUG, emu8086, atau Turbo Debugger.

    \vspace{0.15em}
    \item Catat keadaan awal register dan \eng{flag}.

    \vspace{0.15em}
    \item Jalankan satu instruksi dengan T atau F7.

    \vspace{0.15em}
    \item Catat perubahan register dan \eng{flag}.

    \vspace{0.15em}
    \item Ulangi sampai kesalahan ditemukan.
\end{enumerate}

\vspace{0.15em}
{\footnotesize\color{obeGray}Gunakan P atau F8 untuk melewati prosedur
yang sudah dipercaya.\par}
\Sumber{Bab XVI, Bagian 2}
\end{frame}

% ---------------------------------------------------------
\begin{frame}[fragile]{bug-loop: Gejala dan Penyebab}
\small
Program dimaksudkan menjumlahkan 1 sampai 5, tetapi CX dinaikkan di
dalam perulangan.

\vspace{0.2em}
\begin{kodeasm}[numbers=none]
XOR AX, AX
MOV CX, 5
ULANG:
    ADD AX, CX
    INC CX          ; BUG: CX dinaikkan
    LOOP ULANG
\end{kodeasm}

\vspace{0.15em}
{\footnotesize\color{obeMaroon}LOOP menurunkan CX kembali, sehingga CX
tidak pernah mencapai nol.\par}
\Sumber{Bab XVI, Bagian 2}
\end{frame}

% ---------------------------------------------------------
\begin{frame}{Trace bug-loop}
\footnotesize
\begin{center}
\begin{tabular}{@{}>{\raggedright\arraybackslash}p{1.4cm}
    >{\raggedright\arraybackslash}p{2.6cm}
    >{\raggedright\arraybackslash}p{1.2cm}
    >{\raggedright\arraybackslash}p{1.2cm}
    >{\raggedright\arraybackslash}p{5.0cm}@{}}
    \toprule
    \textbf{Iterasi} & \textbf{Instruksi} & \textbf{AX} &
    \textbf{CX} & \textbf{Keterangan} \\
    \midrule
    1 & ADD AX, CX & 5 & 5 & AX = 0 + 5 \\
    \addlinespace
    1 & INC CX & 5 & 6 & CX naik \\
    \addlinespace
    1 & LOOP ULANG & 5 & 5 & CX turun; lompat \\
    \addlinespace
    2 & ADD AX, CX & 10 & 5 & AX = 5 + 5 \\
    \addlinespace
    2 & INC CX & 10 & 6 & CX naik lagi \\
    \addlinespace
    2 & LOOP ULANG & 10 & 5 & CX tidak pernah nol \\
    \bottomrule
\end{tabular}
\end{center}
\Sumber{Bab XVI, Bagian 2}
\end{frame}

% ---------------------------------------------------------
\begin{frame}[fragile]{Perbaikan bug-loop}
\small
Hapus \code{INC CX}; LOOP sudah menurunkan CX sendiri.

\vspace{0.2em}
\begin{kodeasm}[numbers=none]
XOR AX, AX
MOV CX, 5
ULANG:
    ADD AX, CX
    LOOP ULANG
\end{kodeasm}

\vspace{0.15em}
{\footnotesize\color{obeTeal}Hasil: AX = 15; keluaran
\code{Jumlah 1..5 = 15}.\par}
\Sumber{Bab XVI, Bagian 2}
\end{frame}

% ---------------------------------------------------------
\begin{frame}{Rangkuman Bagian 2}
\small
\begin{itemize}
    \item Single step menjalankan satu instruksi lalu berhenti.

    \vspace{0.15em}
    \item Catat register dan flag setiap langkah.

    \vspace{0.15em}
    \item T/F7 masuk prosedur; P/F8 melewati prosedur.

    \vspace{0.15em}
    \item Perulangan tak berujung terjadi bila pencacah tidak pernah
          mencapai nilai berhenti.

    \vspace{0.15em}
    \item Hapus INC CX memperbaiki bug-loop.
\end{itemize}
\end{frame}

% =========================================================
% BAGIAN 3
% =========================================================
\PemisahBagian{Memeriksa Memori, Stack, dan Memperbaiki Hang}

% ---------------------------------------------------------
\begin{frame}[fragile]{Memeriksa Memori dengan Dump}
\small
Perintah D menampilkan isi memori dalam heksadesimal.

\vspace{0.2em}
\begin{kodeasm}[numbers=none]
D DS:0000           ; awal segmen data
D DS:0010 L 20      ; 20h byte mulai offset 0010h
D SS:SP             ; stack pada puncak saat ini
\end{kodeasm}

\vspace{0.15em}
{\footnotesize\color{obeGray}Baca DS lebih dahulu dengan R. Alamat yang
tampil adalah pasangan segmen:offset, bukan alamat fisik.\par}
\Sumber{Bab XVI, Bagian 3}
\end{frame}

% ---------------------------------------------------------
\begin{frame}[fragile]{bug-stack: Gejala dan Penyebab}
\small
Prosedur mendorong tiga register, tetapi hanya memulihkan dua.

\vspace{0.2em}
\begin{kodeasm}[numbers=none]
HITUNG PROC
    PUSH AX
    PUSH BX
    PUSH CX
    POP CX
    POP BX
    ; POP AX tidak ada
    RET
HITUNG ENDP
\end{kodeasm}

\vspace{0.1em}
{\footnotesize\color{obeMaroon}RET mengambil nilai AX lama, bukan alamat
kembali.\par}
\Sumber{Bab XVI, Bagian 3}
\end{frame}

% ---------------------------------------------------------
\begin{frame}{Trace bug-stack}
\footnotesize
\begin{center}
\begin{tabular}{@{}>{\raggedright\arraybackslash}p{4.0cm}
    >{\raggedright\arraybackslash}p{3.8cm}
    >{\raggedright\arraybackslash}p{4.8cm}@{}}
    \toprule
    \textbf{Tahap} & \textbf{Isi puncak stack} & \textbf{Keterangan} \\
    \midrule
    Sebelum CALL & kosong & stack belum dipakai \\
    \addlinespace
    Sesudah CALL & alamat kembali & disimpan CALL \\
    \addlinespace
    Sesudah tiga PUSH & AX, BX, CX & tiga word tersimpan \\
    \addlinespace
    Sesudah dua POP & AX lama & satu word tersisa \\
    \addlinespace
    RET & AX lama dianggap alamat & lompat ke alamat salah \\
    \bottomrule
\end{tabular}
\end{center}
\Sumber{Bab XVI, Bagian 3}
\end{frame}

% ---------------------------------------------------------
\begin{frame}[fragile]{Perbaikan dan Pemeriksaan Stack Frame}
\begin{columns}[t]
\begin{column}{0.44\textwidth}
{\color{obeTeal}\bfseries Perbaikan}\\[0.15em]
\begin{kodeasm}[numbers=none]
POP CX
POP BX
POP AX          ; perbaikan
RET
\end{kodeasm}

\vspace{0.2em}
{\footnotesize Jumlah PUSH dan POP seimbang.\par}
\end{column}
\begin{column}{0.52\textwidth}
{\color{obeTeal}\bfseries Periksa stack frame}\\[0.15em]
\begin{itemize}\footnotesize
    \item \code{[BP+0]}: BP lama
    \item \code{[BP+2]}: alamat kembali
    \item \code{[BP+4]}: parameter pertama
    \item \code{[BP+6]}: parameter kedua
\end{itemize}

\vspace{0.2em}
{\footnotesize\color{obeGray}Periksa dengan \code{D SS:BP}.\par}
\end{column}
\end{columns}
\Sumber{Bab XVI, Bagian 3}
\end{frame}

% ---------------------------------------------------------
\begin{frame}{Analisis Program yang Berhenti}
\small
\begin{enumerate}
    \item Periksa syarat berhenti perulangan.
    \item Periksa nilai pencacah setiap iterasi.
    \item Periksa keseimbangan PUSH dan POP.
    \item Periksa isi stack pada CALL dan RET.
    \item Periksa apakah program mencapai INT 21h fungsi 4Ch.
\end{enumerate}

\vspace{0.2em}
\begin{alertblock}{Hang belum tentu kerusakan prosesor}
\small Sering kali program hanya terjebak pada perulangan yang
syaratnya tidak pernah terpenuhi.
\end{alertblock}
\Sumber{Bab XVI, Bagian 3}
\end{frame}

% ---------------------------------------------------------
\begin{frame}{Rangkuman Bagian 3}
\small
\begin{itemize}
    \item Perintah D menampilkan memori atau stack dalam
          heksadesimal.

    \vspace{0.15em}
    \item Alamat DEBUG adalah segmen:offset, bukan alamat fisik.

    \vspace{0.15em}
    \item Stack tidak seimbang membuat RET mengambil alamat salah.

    \vspace{0.15em}
    \item Perbaikan: tambahkan POP AX sebelum RET.

    \vspace{0.15em}
    \item Analisis hang: perulangan, pencacah, stack, dan instruksi
          keluar.
\end{itemize}
\end{frame}

% =========================================================
% CONTOH PROGRAM
% =========================================================
\begin{frame}[fragile]{Contoh 1: Perulangan Tak Berujung}
\begin{columns}[t]
\begin{column}{0.47\textwidth}
{\color{obeMaroon}\bfseries Sebelum perbaikan}\\[0.15em]
\begin{kodeasm}[numbers=none]
XOR AX, AX
MOV CX, 5
ULANG:
    ADD AX, CX
    INC CX
    LOOP ULANG
\end{kodeasm}
\end{column}
\begin{column}{0.49\textwidth}
{\color{obeTeal}\bfseries Sesudah perbaikan}\\[0.15em]
\begin{kodeasm}[numbers=none]
XOR AX, AX
MOV CX, 5
ULANG:
    ADD AX, CX
    LOOP ULANG
\end{kodeasm}
\end{column}
\end{columns}

\vspace{0.2em}
{\footnotesize\color{obeTeal}Keluaran setelah perbaikan:
\code{Jumlah 1..5 = 15}.\par}
\Sumber{Bab XVI, Contoh 1}
\end{frame}

% ---------------------------------------------------------
\begin{frame}[fragile]{Contoh 2: Stack Tidak Seimbang}
\begin{columns}[t]
\begin{column}{0.47\textwidth}
{\color{obeMaroon}\bfseries Sebelum perbaikan}\\[0.15em]
\begin{kodeasm}[numbers=none]
PUSH AX
PUSH BX
PUSH CX
POP CX
POP BX
RET
\end{kodeasm}
\end{column}
\begin{column}{0.49\textwidth}
{\color{obeTeal}\bfseries Sesudah perbaikan}\\[0.15em]
\begin{kodeasm}[numbers=none]
PUSH AX
PUSH BX
PUSH CX
POP CX
POP BX
POP AX
RET
\end{kodeasm}
\end{column}
\end{columns}

\vspace{0.2em}
{\footnotesize\color{obeTeal}Setelah perbaikan, program kembali ke
pemanggil dan mencetak \code{Selesai.}\par}
\Sumber{Bab XVI, Contoh 2}
\end{frame}

% =========================================================
% PRAKTIKUM
% =========================================================
\begin{frame}{Praktikum: Menemukan dan Memperbaiki Kesalahan}
\begin{block}{Judul praktikum}
\small Menemukan dan memperbaiki kesalahan dengan \eng{single step}
dan \eng{dump}\\[0.15em]
\textbf{Sub-CPMK:} 4.1--4.3 \quad\textbullet\quad
\textbf{Waktu:} 100 menit \quad\textbullet\quad
\textbf{Bentuk:} individu
\end{block}

\vspace{0.2em}
\begin{columns}[t]
\begin{column}{0.46\textwidth}
{\color{obeBlue}\bfseries Alat dan bahan}
\begin{itemize}\footnotesize
    \item Komputer dengan emu8086
    \item DOSBox dengan TASM, TLINK, dan DEBUG
    \item Empat berkas \file{bug-*.asm}
    \item Lembar laporan debugging
\end{itemize}
\end{column}
\begin{column}{0.50\textwidth}
{\color{obeBlue}\bfseries Langkah kerja}
\begin{enumerate}\footnotesize
    \item Reproduksi hang pada \file{bug-loop.asm}.
    \item Single step; catat AX dan CX.
    \item Hapus INC CX; uji ulang.
    \item Reproduksi hang pada \file{bug-stack.asm}.
    \item Periksa stack; tambahkan POP AX.
\end{enumerate}
\end{column}
\end{columns}

\vspace{0.2em}
{\footnotesize\textbf{Dikumpulkan:} dua berkas perbaikan, tabel
penelusuran AX/CX/SP/stack, tangkapan layar, dan laporan
debugging.\par}
\Sumber{Bab XVI, Praktikum}
\end{frame}

% =========================================================
% LATIHAN
% =========================================================
\begin{frame}{Latihan}
\small
\begin{enumerate}
    \item Jelaskan perbedaan kesalahan sintaks, logika, dan hang;
          berikan satu contoh masing-masing.

    \vspace{0.15em}
    \item Sebutkan fungsi R, T, P, D, dan G pada DEBUG. Kapan P lebih
          tepat daripada T?

    \vspace{0.15em}
    \item Jelaskan mengapa CX tidak pernah nol pada
          \file{bug-loop.asm} dan tentukan perbaikannya.

    \vspace{0.15em}
    \item Jelaskan mengapa RET melompat ke alamat salah pada
          \file{bug-stack.asm} dan tentukan perbaikannya.

    \vspace{0.15em}
    \item Sebutkan langkah pemeriksaan untuk menemukan penyebab
          program yang berhenti.
\end{enumerate}
\end{frame}

% =========================================================
% RANGKUMAN
% =========================================================
\begin{frame}{Rangkuman}
\small
\begin{itemize}
    \item Kesalahan dibagi menjadi sintaks, logika, dan hang.

    \vspace{0.15em}
    \item DEBUG menyediakan R, T, P, D, dan G untuk penelusuran.

    \vspace{0.15em}
    \item Penelusuran mencatat register dan flag setiap instruksi.

    \vspace{0.15em}
    \item Perulangan tak berujung: pencacah tidak mencapai nilai
          berhenti.

    \vspace{0.15em}
    \item Stack tidak seimbang: jumlah PUSH dan POP berbeda.

    \vspace{0.15em}
    \item Perbaikan harus diverifikasi dengan trace ulang.
\end{itemize}
\end{frame}

% ---------------------------------------------------------
\begin{frame}{Cek Pemahaman}
\footnotesize
\begin{enumerate}
    \item Perintah DEBUG untuk menjalankan satu instruksi adalah
          \dots
          \begin{enumerate}
              \item R \quad \item T \quad \item D \quad \item G
          \end{enumerate}

    \vspace{0.2em}
    \item Perintah yang menjalankan satu instruksi tetapi melewati
          prosedur adalah \dots
          \begin{enumerate}
              \item T \quad \item P \quad \item D \quad \item R
          \end{enumerate}

    \vspace{0.2em}
    \item Penyebab hang pada \file{bug-stack.asm} adalah \dots
          \begin{enumerate}
              \item terlalu banyak POP
              \item satu POP tidak ada
              \item CALL tidak menyimpan alamat kembali
              \item prosedur tidak memiliki RET
          \end{enumerate}
\end{enumerate}
\end{frame}

% =========================================================
% PENUTUP
% =========================================================
\begin{frame}{Penutup}
\begin{block}{Menuju UAS dan Proyek Akhir}
\small
\begin{itemize}
    \item Selesaikan praktikum: dua berkas perbaikan, tabel
          penelusuran, tangkapan layar, dan laporan debugging.
    \item Ulangi Bab X--XVI untuk UAS: aritmatika, logika, mode
          pengalamatan, kontrol, prosedur/macro, array/string, dan
          debugging.
    \item Tidak ada deck materi baru pada Minggu 16 (UAS).
\end{itemize}
\end{block}

\vspace{0.4em}
\begin{center}
{\color{obeBlue}\sffamily\large\bfseries Terima kasih\par}
\vspace{0.15em}
{\color{obeGray}\sffamily\small \NamaMK\ (\KodeMK)
\textbullet\ \SemesterTeks\par}
\end{center}
\end{frame}

\end{document}
